\begin{essaywithapp}

	Выпускная квалификационная работа посвящена разработке и экспериментальной проверке модуля мультиобъектного отслеживания (MOT) для роботизированного транспортного средства.

	\textbf{\textit{Объект исследования ---}} интеллектуальные системы помощи водителю (ИСПВ), а также алгоритмические и аппаратные компоненты автономных транспортных средств.

	\textbf{\textit{Цель работы ---}} разработка системы мультиобъектного отслеживания для роботизированного ТС на основе нелинейного алгоритма предсказания, обеспечивающая устойчивое определение положения объектов в условиях их частичного перекрытия и нагруженной дорожной ситуации.

	\textbf{\textit{Методы исследования ---}} анализ, классификация, сравнение, моделирование, тестирование.

	\textbf{\textit{Результаты и научная новизна ---}} реализована система мультобъектного отслеживания на основе нелинейного фильтра предсказания,позволяющий предсказывать положение объекта при их частичном или полном перекрытии в условиях нагруженной дорожной ситуации; проведено тестирование и анализ метрик, а также сравнение с аналогичными системами МОТ; разработан алгоритм поиска наиболее важного объекта.

	\textbf{\textit{Область применения ---}} автономные роботизированные ТС, системы интеллектуальной помощи водителю (ИСПВ).

	\textbf{\textit{Рекомендации по внедрению ---}} использование системы в общественном и коммерческом транспорте для повышения безопасности на дорогах общего пользования. 

\end{essaywithapp}